\documentclass[10pt,a4paper]{amsart}
\usepackage[margin=19mm]{geometry}
\usepackage{amsmath,amssymb,mathtools}
\usepackage[scr=rsfs,cal=euler]{mathalfa}
\usepackage{xfrac}
\usepackage[unicode,hidelinks,bookmarks=false]{hyperref}
\newcommand{\ie}{i.e.\ }
\newcommand{\cf}{cf.\ }
\newcommand{\resp}{respectively\ }
\newcommand{\Subsection}{\subsection}
\providecommand{\nobreakdash}{\nobreak\hbox{-}}
\setlength{\emergencystretch}{2em}
\multlinegap0pt
\usepackage{bm}
\usepackage{bbm}
\usepackage{dsfont}
\usepackage{xspace}
\usepackage{cancel}
\usepackage{mathtools}
\usepackage{amsthm}
\makeatletter
\@ifundefined{theorem}{\newtheorem{theorem}{Theorem}}{}
\@ifundefined{corollary}{\newtheorem{corollary}[theorem]{Corollary}}{}
\@ifundefined{proposition}{\newtheorem{proposition}[theorem]{Proposition}}{}
\makeatother
\def\p{\varphi}
\title[Geometry of dyadic polygons II: isomorphisms of dyadic trian]{Geometry of dyadic polygons II: isomorphisms of dyadic triangles}
\author{A. Mu\'cka, A.B. Romanowska}
\date{Source: arXiv:2510.07244v1 (CC BY 4.0)}
\begin{document}
\maketitle

\noindent\emph{Source and licence.} Geometry of dyadic polygons II: isomorphisms of dyadic triangles. Version arXiv:2510.07244v1 (CC BY 4.0), distributed under the Creative Commons Attribution 4.0 International licence, as recorded in the arXiv licence metadata for this exact version and shown on the abstract page https://arxiv.org/abs/2510.07244v1. Full rights notice: licenses/arxiv-2510.07244-CC-BY-4.0.txt.\par\smallskip

\noindent\textit{Editorial scope.} Counted unit: the Proposition whose source label is \texttt{P:mmm} in the article (source0.tex lines 789--799, source label \texttt{P:mmm}), delivered together with the article's own complete original proof of it (source0.tex lines 800--854, one \texttt{proof} environment of 55 lines). No mathematics is written, repaired or re-derived: statement and proof are the article's own text line for line. Required context retained uncounted: C:mmm (lines 760--767). Every definition, display and label of those lines is retained unchanged. Omitted from this package and not redistributed: 1--759, 768--788, 855--1172. Nothing inside a retained excerpt is omitted or altered: every retained body is delivered exactly as the article's lines are written, so no omission receipt of any kind accompanies this package. Citations: the retained excerpts carry no citation command at all, so no bibliography entry of the article is transcribed and no reference is redistributed. Reference adaptation: none was needed and none was made; every reference inside the delivered unit resolves inside it, so no unresolved reference or citation is produced. Printed slips kept literal: no printed word, symbol, hypothesis or number of the counted statement or of its proof was altered. Layout and class: the article's own class is \texttt{amsart} with packages including amssymb, latexsym, float, hyperref, color; the wrapper is the lane's pinned \texttt{amsart} editorial wrapper at 10pt on A4, which restates the theorem-like environments the retained text needs and adds the article's own macro definitions that the retained text needs (\texttt{\textbackslash p}), with extra wrapper packages (bm, bbm, dsfont, xspace, cancel, mathtools). The delivered numbering is the unit's own. Assets: no retained span contains a figure environment, a graphics-inclusion command, a PDF-page inclusion, a table or a TikZ picture; no image file is copied into this package, the package declares no asset dependency and no image file is redistributed with it. Licence: version arXiv:2510.07244v1 of this article is distributed under the Creative Commons Attribution 4.0 International licence, as recorded in the arXiv licence metadata for this exact version and shown on the abstract page https://arxiv.org/abs/2510.07244v1. A full rights notice is delivered at licenses/arxiv-2510.07244-CC-BY-4.0.txt. Characters: every retained excerpt is pure ASCII, so the delivered engine, XeLaTeX with its default Latin Modern text font, reports no missing character.
\begin{corollary}\label{C:mmm}
Let $T = ABC$ be a representative hat $T_{i,j,m}$ with $i \neq m$, and with boundary type
$(m,m,m)$. Then
\begin{enumerate}
\item[(a)] $i = km$ and $j = lm$ for some relatively prime odd integers $k$ and $l$;
\item[(b)] $m - i = nm$ for some even integer $n$ co-prime to $l$.
\end{enumerate}
\end{corollary}

\begin{proposition}\label{P:mmm}
Let $T = ABC$ be a representative hat $T_{i,j,m}$ of boundary type $(m,m,m)$. Let $\p: T
\rightarrow T$ be a mapping. Then the following statements hold.
\begin{enumerate}
\item[(a)] The mapping $\p$ is an automorphism of the hat $T$ taking either $A$ to $B$, $B$ to $C$ and $C$ to $A$ or $A$ to $C$, $B$ to $A$ and $C$ to $B$ precisely when $l$ divides $k^2 - k + 1$, with $k$ and $l$ defined as in Corollary~\ref{C:mmm}.
\item[(b)] The automorphism group of $T$ contains a subgroup isomorphic to the cyclic group
    $C_3$ if and only if $l$ divides $k^2 - k + 1$.
\item[(c)] If $T$ has no automorphisms fixing precisely one vertex, then the automorphism
    group of $T$ is isomorphic to the cyclic group $C_3$ if and only if $l$ divides $k^2 - k + 1$.
\end{enumerate}
\end{proposition}

\begin{proof}
(a) First note that by Corollary~\ref{C:mmm}, $i = km$ and $j = lm$ for some relatively prime
odd integers $k$ and $l$. Then translate the hat $T$ to the isomorphic hat $T' = A'B'C'$ with
$A' = (-m,0)$, $B' = (i-m,j)$ and $C' = (0,0)$. The following matrix $M$ provides an
isomorphism from the hat $T'$ onto the hat $T$, taking $A'$ to $B$, $B'$ to $C$, and $C'$ to
$A$.
\begin{align*}
M &=
	\begin{bmatrix}
		i-m&j\\
		-m&0
	\end{bmatrix}^{-1}
	\begin{bmatrix}
		m&0\\
		i&j
	\end{bmatrix}
    =\begin{bmatrix}
	-i/m&-j/m\\
	(m^2-im+i^2)/mj&(-m+i)/m
    \end{bmatrix}\\
    &= \begin{bmatrix}
     -k&-l\\
	(1-k+k^2)/l&k-1
    \end{bmatrix}.
\end{align*}
Note that $\det M = 1$, and that the matrix $M$ is dyadic precisely when $l$ is an integer dividing
$k^2 - k + 1$.

Similarly, the matrix $M'$ below provides an isomorphism from the hat $T'$ onto the hat $T$,
taking $A'$ to $C$, $B'$ to $A$, and $C'$ to $B$.
\begin{align*}
M' &=
	\begin{bmatrix}
		-i&-j\\
		m-i&-j
	\end{bmatrix}^{-1}
	\begin{bmatrix}
		m&0\\
		i&j
	\end{bmatrix}
    =\begin{bmatrix}
	i-m/m&j/m\\
	(mi-m^2 -i^2)/mj&-i/m
    \end{bmatrix}\\
    &= \begin{bmatrix}
     k-1&l\\
	(-1+k-k^2)/l&-k
    \end{bmatrix}.
\end{align*}
The matrix $M'$ is dyadic precisely when $l$ is an integer dividing $k^2 - k + 1$.

Let $\p$ be the composition of the first translation and the isomorphism given by the matrix $M$ or $M'$. It is clear that $\p$ is an automorphism satisfying the required conditions precisely when $l$ divides $k^2 - k + 1$.

Now (b) follows directly by (a), since each version of the automorphism $\p$ generate the other and the identity map. Then (c) is a consequence of (b).
\end{proof}

\end{document}
